\documentclass[10pt,a4paper]{amsart}
\usepackage[margin=19mm]{geometry}
\usepackage{amsmath,amssymb,mathtools}
\usepackage[scr=rsfs,cal=euler]{mathalfa}
\usepackage{xfrac}
\usepackage[unicode,hidelinks,bookmarks=false]{hyperref}
\newcommand{\ie}{i.e.\ }
\newcommand{\cf}{cf.\ }
\newcommand{\resp}{respectively\ }
\newcommand{\Subsection}{\subsection}
\providecommand{\nobreakdash}{\nobreak\hbox{-}}
\setlength{\emergencystretch}{2em}
\multlinegap0pt
\usepackage{bm}
\usepackage{bbm}
\usepackage{dsfont}
\usepackage{xspace}
\usepackage{cancel}
\usepackage{mathtools}
\usepackage{amsthm}
\makeatletter
\@ifundefined{theorem}{\newtheorem{theorem}{Theorem}}{}
\@ifundefined{lemma}{\newtheorem{lemma}[theorem]{Lemma}}{}
\makeatother
\DeclareMathOperator{\dimqa}{dim_{qA}}
\newcommand{\N}{\mathbb{N}}
\newcommand{\bu}{\mathbf{u}}
\newcommand{\ccc}{\mathbf{c}}
\newcommand{\de}{\delta}
\newcommand{\eps}{\varepsilon}
\newcommand{\f}{\frac}
\newcommand{\iii}{\mathbf{i}}
\newcommand{\jjj}{\mathbf{j}}
\newcommand{\kk}{\mathcal{K}}
\newcommand{\ok}{\overline{\mathcal{K}}}
\newcommand{\om}{\overline{M}}
\newcommand{\uc}{\underline{c}}
\newcommand{\uk}{\underline{\mathcal{K}}}
\newcommand{\um}{\underline{M}}
\renewcommand{\geq}{\geqslant}
\renewcommand{\ge}{\geqslant}
\renewcommand{\leq}{\leqslant}
\renewcommand{\le}{\leqslant}
\title[Assouad and quasi-Assouad dimensions of Moran sets]{Assouad and quasi-Assouad dimensions of Moran sets}
\author{Junjie Miao, Minghui Xu}
\date{Source: arXiv:2511.09255v1 (CC BY 4.0)}
\begin{document}
\maketitle

\noindent\emph{Source and licence.} Assouad and quasi-Assouad dimensions of Moran sets. Version arXiv:2511.09255v1 (CC BY 4.0), distributed under the Creative Commons Attribution 4.0 International licence, as recorded in the arXiv licence metadata for this exact version and shown on the abstract page https://arxiv.org/abs/2511.09255v1. Full rights notice: licenses/arxiv-2511.09255-CC-BY-4.0.txt.\par\smallskip

\noindent\textit{Editorial scope.} Counted unit: the Theorem whose source label is \texttt{xiaoyu} in the article (source0.tex lines 1042--1047, source label \texttt{xiaoyu}), delivered together with the article's own complete original proof of it (source0.tex lines 1048--1175, one \texttt{proof} environment of 128 lines). No mathematics is written, repaired or re-derived: statement and proof are the article's own text line for line. Required context retained uncounted: club (lines 405--407); def\_t*tt* (lines 498--504); compensate (lines 875--880); numberofcap (lines 1003--1012). Every definition, display and label of those lines is retained unchanged. Omitted from this package and not redistributed: 1--404, 408--497, 505--874, 881--1002, 1013--1041, 1176--2148. Nothing inside a retained excerpt is omitted or altered: every retained body is delivered exactly as the article's lines are written, so no omission receipt of any kind accompanies this package. Citations: the retained excerpts carry no citation command at all, so no bibliography entry of the article is transcribed and no reference is redistributed. Reference adaptation: none was needed and none was made; every reference inside the delivered unit resolves inside it, so no unresolved reference or citation is produced. Printed slips kept literal: no printed word, symbol, hypothesis or number of the counted statement or of its proof was altered. Layout and class: the article's own class is \texttt{amsart} with packages including amsmath, amssymb, graphicx, IEEEtrantools, enumerate, tikz, comment, bm, mathtools, microtype, geometry, xcolor, hyperref, fancyhdr; the wrapper is the lane's pinned \texttt{amsart} editorial wrapper at 10pt on A4, which restates the theorem-like environments the retained text needs and adds the article's own macro definitions that the retained text needs (\texttt{\textbackslash N}, \texttt{\textbackslash bu}, \texttt{\textbackslash ccc}, \texttt{\textbackslash de}, \texttt{\textbackslash dimqa}, \texttt{\textbackslash eps}, \texttt{\textbackslash f}, \texttt{\textbackslash ge}, \texttt{\textbackslash geq}, \texttt{\textbackslash iii}, \texttt{\textbackslash jjj}, \texttt{\textbackslash kk}, \texttt{\textbackslash le}, \texttt{\textbackslash leq}, \texttt{\textbackslash ok}, \texttt{\textbackslash om}, \texttt{\textbackslash uc}, \texttt{\textbackslash uk}, \texttt{\textbackslash um}), with extra wrapper packages (bm, bbm, dsfont, xspace, cancel, mathtools). The delivered numbering is the unit's own. Assets: no retained span contains a figure environment, a graphics-inclusion command, a PDF-page inclusion, a table or a TikZ picture; no image file is copied into this package, the package declares no asset dependency and no image file is redistributed with it. Licence: version arXiv:2511.09255v1 of this article is distributed under the Creative Commons Attribution 4.0 International licence, as recorded in the arXiv licence metadata for this exact version and shown on the abstract page https://arxiv.org/abs/2511.09255v1. A full rights notice is delivered at licenses/arxiv-2511.09255-CC-BY-4.0.txt. Characters: every retained excerpt is pure ASCII, so the delivered engine, XeLaTeX with its default Latin Modern text font, reports no missing character.
	\begin{equation}\label{club}
		\lim_{k\to\infty}\frac{\log \uc_k}{\log \om_k}=0.
	\end{equation}

	\begin{equation}\label{def_t*tt*}
		\begin{split}
			&t^*=\lim_{\eta\to 0}\limsup_{l\to\infty} \sup_{k\in \ok_{l+1,\eta}} s_{k+1,k+l}, \qquad
			t=\lim_{\eta\to 0}\limsup_{l\to\infty}\sup_{k\in \kk_{l,\eta}} s_{k+1,k+l}, \\
			&t_*=\lim_{\eta\to 0}\limsup_{l\to\infty}\sup_{k\in \uk_{l,\eta}} s_{k+1,k+l}, 
		\end{split}
	\end{equation}

	\begin{lemma}\label{compensate}
	Let $\tilde{D}\subset D^*$ be such that $\{J_{\bu}:\bu\in\tilde{D}\}$ forms a finite, non-overlapping covering of the Moran set $E$. Then for any $s>0$, there exist $k_1,k_2$ with $\min_{\bu\in\tilde{D}}|\bu|\leq k_1,k_2\leq \max_{\bu\in\tilde{D}}|\bu|$ such that
		\[
		\sum_{\bu\in{D}_{k_1}}c^{s}_{\bu}\leq\sum_{\bu\in\tilde{D}}c^{s}_{\bu}\leq\sum_{\bu\in{D}_{k_2}}c^{s}_{\bu}.
		\]
	\end{lemma}

	\begin{lemma}\label{numberofcap}
Suppose  that $\mathcal{M}(J,\{n_k\},\{\ccc_k\})$ satisfies \eqref{club}.  Then, for all $E\in\mathcal{M}(J,\{n_k\},\{\ccc_k\})$,  the limit
		\[
		\lim_{\de\to 0}\frac{\log \sup_{x\in E}\#\{\bu\in D_{\iii}(\de):B(x,c_{\iii}\de)\cap J_{\bu}\neq \emptyset\}}{-\log c_{\iii}\de}=0
		\]
		holds uniformly in  $\iii\in D^*$. In particular, if $\iii$ is the empty word, we have 
	\begin{equation}\label{AWSC}
		\lim_{\de\to 0}\frac{\log \sup_{x\in E}\#\{\bu\in D(\de):B(x,\de)\cap J_{\bu}\neq \emptyset\}}{-\log\de}=0.
		\end{equation}
	\end{lemma}

	\begin{theorem}\label{xiaoyu}
		Let $ \mathcal{M}(J,\{n_k\},\{\ccc_k\})$ satisfy \eqref{club}. Then every  $E\in\mathcal{M}(J,\{n_k\},\{\ccc_k\})$ satisfies
		\[
		\dimqa E\le t.
		\]
	\end{theorem}

	\begin{proof}%[Proof of Theorem \ref{xiaoyu}]
		Fix $\eta>0$, and  set  $s^{*}(\eta)=\limsup_{l\to\infty}\sup_{k\in \kk_{l,\eta}} s_{k+1,k+l}$. By \eqref{def_t*tt*}, it is clear that $\lim_{\eta \to 0} s^{*}(\eta)=t^*$,  and it suffices to prove that
		\[
		h_E(2\eta)\le s^*(\eta).
		\]
		
		Arbitrarily choose $\eps>0$.	Since \eqref{club} holds, it folows from
		\[
		\f{\f{\log \uc_{k+1}}{\log \om_k}}{1+\f{\log \uc_{k+1}}{\log \om_k}}\le\f{\log \uc_{k+1}}{\log \om_{k+1}}
		\]
		 that $\lim_{k\to\infty}\f{\log \uc_{k+1}}{\log \om_k}=0$. Hence, by induction, we also have $\lim_{k\to\infty}\f{\log \uc_{k+l}}{\log \om_k}=0$ for any $l$.   Note that  $\sum
		_{j=1}^{n_{k}}c_{k,j}^{\,d}\leq 1$ for all $k>0$, it follows that
		\[
		\f{\log n_{k+1}\cdots n_{k+l}}{-\log \om_{k-1}}\le\f{d\log \uc_{k+1}\cdots \uc_{k+l}}{\log \om_{k-1}} .
		\]
Combining these with  Lemma \ref{numberofcap},
		there exist $k_0, k'_0 \in {\N}^+,$ and $\de_0\in (0,\eta)$ such that
		\begin{eqnarray}
			\sup_{k\in \kk_{l,\eta}} s_{k+1,k+l}<s^*(\eta)+2\eps, \qquad &&l\ge k_0,\label{sk}   \\
			\f{\log \uc_k}{\log \om_k}<\frac{1}{2},\qquad  \frac{\om_{|\iii|+k}^{\eps}}{\uc_{|\iii|}^{s^*(\eta)+2\eps}\uc_{|\iii|+k}^{s^*(\eta)+2\eps}}<1,\qquad && k\ge k_0,\ \iii\in D^*,\vspace{0.5em}\label{dkmk}\\
			n_{k+1}n_{k+2}\cdots n_{k+k_0-1}<\om_{k-1}^{- {\eps}},\qquad &&k\ge k'_0,\label{nk}\vspace{0.2em}\\
			\sup_{x\in E}\#\{\bu\in D(\delta):B(x,\de)\cap J_{\bu}\neq \emptyset\}<\de^{-{\eps}},\qquad &&0<\de<\de_0.\label{numberofcut'}
		\end{eqnarray}
		Given $\iii\in D^*$, for each $k\in{\N}^+$ and $\delta>0$, we write
		\[
		D_{\iii}(\de,k)=\{\iii*\jjj\in D^*:c_{\jjj}\le \de<c_{\jjj^-},|\jjj|=k\}
		\]
		and $D(\de,k)=D_{\emptyset}(\de,k)$. For each $\bu=\iii*\jjj\in D_{\iii}(\de,k)$, it is clear that
		$$
		\de^{s^*(\eta)+2\eps}< c_{\jjj^-}^{s^*(\eta)+2\eps} \leq \Big(\frac{c_{\bu}}{c_{\iii}\uc_{|\iii|+k}}\Big)^{s^*(\eta)+2\eps} \le  \frac{c_{\bu}^{s^*(\eta)+\eps}\om_{|\iii|+k}^{\eps}}{c_{\iii^-}^{s^*(\eta)+2\eps}\uc_{|\iii|}^{s^*(\eta)+2\eps}\uc_{|\iii|+k}^{s^*(\eta)+2\eps}}.
		$$
		For $k\ge k_0$,  by \eqref{dkmk}, it follows that
		\begin{equation}\label{estdelta1}
			\begin{aligned}
				\de^{s^*(\eta)+2\eps}&< {c_{\iii^-}^{-(s^*(\eta)+2\eps)}}  c_{\bu}^{s^*(\eta)+\eps} .
			\end{aligned}
		\end{equation}
		
		
		Given $r,R$ satisfying $0<r<R^{1+2\eta}<R<\min \{ \de_0,\om_{k_0},\om_{k'_0}\}$,  for every $\iii\in D(R)$  and $\bu\in D_\iii\left(\f rR\right)$, we have $|\iii|\ge \max\{k_0,k'_0\}$ and $\frac{c_\bu}{c_\iii}\leq \frac{r}{R}$. By \eqref{dkmk}, we have
		$$
		\log \om_{|\iii|}\geq 2(\log \om_{|\iii|}-\log\uc_{|\iii|})\geq 2(\log c_{\iii}-\log\uc_{|\iii|}) \geq 2\log c_{\iii^-} \geq 2\log R.
		$$
		Immediately, it follows that 
		\begin{align*}
			\f{\log \um_{|\iii|+1,|\bu|}}{\log \om_{|\iii|}}
			&\ge \f{\log{c_\bu}-\log{c_\iii}}{\log \om_{|\iii|}} >\f{\log\f{r}R}{2\log R} >\eta, 
		\end{align*}
		and this implies  $|\iii|\in \kk_{|\bu|-|\iii|,\eta}$.
		
		
		For every $\iii\in D(R)$, define $k_1=\max\big\{k:D_\iii(\f rR,k)\neq\emptyset\big\} $. Then we have the following two cases.
		
		\noindent	Case A: If $k_1 <k_0$,  by \eqref{nk}, we have
		\begin{align*}
			\# D_{\iii}\Big(\f rR\Big)
			&=\sum_{k=1}^{k_1}\# D_{\iii}\Big(\f rR,k\Big)\ \le n_{|\iii|+1}n_{|\iii|+2}\cdots n_{|\iii|+k_0-1} <\om_{|\iii|-1}^{-{\eps}} <c_{\iii^-}^{-{\eps}}.
		\end{align*}
		Case B: If $k_1 \ge k_0$, set $	k'_1=\min\big\{k: D_\iii(\f rR,k)\neq\emptyset,k\ge k_0\big\}$, and  we obtain
		\begin{equation}\label{numberofcut}
			\# D_{\iii}\Big(\f rR\Big)
			=\sum_{k=1}^{k_0-1}\# D_{\iii}\Big(\f rR,k\Big)+\sum_{k=k'_1}^{k_1}\# D_{\iii}\Big(\f rR,k\Big)<c_{\iii^-}^{-{\eps}}+\sum_{k=k'_1}^{k_1}\#D_{\iii}\Big(\f rR,k\Big).
		\end{equation}
	Let
		\begin{align*}
			Q=\Big(\left\{J_\bu:\bu\in D_\iii\Big(\f rR,k'_1\Big)\right\}\Big\backslash \Big\{J_{\bu}:\bu=\mathbf{\tau}|_{|\iii|+k'_1},\mathbf{\tau}\in\bigcup_{k=k'_1}^{k=k_1}D_{\iii}\Big(\f rR,k\Big)\Big\}\Big)\\
			\bigcup \Big\{J_{\bu}:\bu\in \bigcup_{k=k'_1}^{k=k_1}D_{\iii}\Big(\f rR,k\Big)\Big\}.
		\end{align*}
	Then $Q$ is a finite, non-overlapping covering of $J_\iii \cap E$, and by \eqref{estdelta1}, we have 
		$$	
		\sum_{k=k'_1}^{k_1}\# D_{\iii}\Big(\f rR,k\Big)\Big(\f rR\Big)^{s^*(\eta)+2\eps} <\frac1 {c_{\iii^-}^{s^*(\eta)+2\eps}}\sum_{k=k'_1}^{k_1}\sum_{\bu\in D_{\iii}(\f rR,k)}c_{\bu}^{s^*(\eta)+2\eps} 
		\le\frac1 {c_{\iii^-}^{s^*(\eta)+2\eps}}\sum_{J_\bu\in Q}c_{\bu}^{s^*(\eta)+2\eps}
		$$	
		By Lemma \ref{compensate},  there exists  $k_1'\leq k^*\leq k_1$ such that  
		$$
		\sum_{J_\bu\in Q}c_{\bu}^{s^*(\eta)+2\eps}\leq c_{\iii}^{s^*(\eta)+2\eps}\sum_{\jjj\in  D_{|\iii|+1,|\iii|+k^*}}c_{\jjj}^{s^*(\eta)+2\eps}= c_{\iii}^{s^*(\eta)+2\eps} \prod_{k=|\iii|+1}^{|\iii|+k^*}\sum_{j=1}^{n_k}c_{k,j}^{s^*(\eta)+2\eps}.
		$$
		Note that $|\iii|\in \kk_{k'_1,\eta}\subseteq \kk_{k^*,\eta}$.   Combining these with	\eqref{sk}, we have that
		\begin{equation*} %\label{estdelta2}
			\begin{aligned}
				\sum_{k=k'_1}^{k_1}\# D_{\iii}\Big(\f rR,k\Big)\Big(\f rR\Big)^{s^*(\eta)+2\eps}
				&<c_{\iii^-}^{-\eps}\prod_{k=|\iii|+1}^{|\iii|+k^*}\sum_{j=1}^{n_k}c_{k,j}^{s^*(\eta)+2\eps}<c_{\iii^-}^{-\eps}.
			\end{aligned}
		\end{equation*}  
		By  \eqref{numberofcut}, it follows that
		\begin{equation*}
			\#D_{\iii}\Big(\f rR\Big)< c_{\iii^-}^{-{\eps}}+c_{\iii^-}^{-{\eps}}\Big(\f Rr \Big)^{s^*(\eta)+2\eps}< 2c_{\iii^-}^{-{\eps}}\Big(\f Rr\Big)^{s^*(\eta)+2\eps}.
		\end{equation*}
		
		Combining Case A and Case B, we obtain that
		\begin{equation}\label{numberofcut2}
			\#D_{\iii}\left(\f rR\right)\le 2c_{\iii^-}^{-{\eps}}\Big(\f Rr\Big)^{s^*(\eta)+2\eps} \le 2R^{-{\eps}}\Big(\f Rr\Big)^{s^*(\eta)+2\eps},
		\end{equation}
		for $0<r<R^{1+2\eta}<R<\min \{ \de_0,\om_{k_0},\om_{k'_0} \}$ and $\iii\in D(R)$.
		\iffalse
		Obviously, $\# D(\de)\de^{s^*(\eta)+\eps}<\# D(\de)\le \# D_{k_0}\le N^{k_0}$ when $\om_{k_0}\le\de<1$. It follows that
		\begin{equation}\label{numberofcutset1}
			\# D_{\iii}(\de)=\# D(\de)<N^{k_0}\de^{-s^*(\eta)-\eps}
		\end{equation}
		for all $\iii\in D$ and all $0<\de<1$.
		\fi 
		
		Fix $x\in E$. Next we estimate $N_r(B(x,R)\cap E)$ for $0<r<R^{1+2\eta}<R<\min\left\{ \de_0,\om_{k_0},\om_{k'_0}\right\}$. It is clear that
		\[
		B(x,R)\cap E\subset \bigcup_{\iii\in D(R),J_{\iii}\cap B(x,R)\cap E\neq \emptyset} J_{\iii}\cap E.
		\]
		For each $\iii\in D(R)$ with $J_{\iii}\cap B(x,R)\cap E\neq \emptyset$, we have
		\[
		J_{\iii}\cap E\subseteq\bigcup_{\bu\in D_{\iii}\left(\frac r R\right)} J_{\bu}.
		\]
		For each $\bu\in D_{\iii}\left(\frac r R\right)$, choose a point $x_\bu\in J_\bu\cap E$. Then $	J_{\bu}\subseteq B(x_{\bu},r) $ since  $c_{\bu}\le  r$ and $|J|=1$.
		Hence, it follows that 
		\[
		B(x,R)\cap E\subset\bigcup_{\iii\in D(R),J_{\iii}\cap B(x,R)\cap E\neq \emptyset}\ \bigcup_{\bu\in D_{\iii}\left(\frac r R\right)}B(x_{\bu},r).
		\]
		Combining it with \eqref{numberofcut2}, we have 
		\begin{align*}
			N_r(B(x,R)\cap E)&\le \sum_{\substack{\iii\in D(R)\\J_{\iii}\cap B(x,R)\cap E\neq \emptyset}}\# D_{\iii}\left(\frac r R\right)<2R^{-{\eps}}\sum_{\substack{\iii\in D(R)\\J_{\iii}\cap B(x,R)\cap E\neq \emptyset}}\Big(\frac R r\Big)^{s^*(\eta)+2\eps} 
		\end{align*}	
		Since $R^{-\eta}<\f Rr$, by \eqref{numberofcut'}, it follows that 
		\begin{align*}
			N_r(B(x,R)\cap E)
			&<2R^{-{\eps}}\Big(\frac R r\Big)^{s^*(\eta)+2\eps}\sup_{x\in E}\#\{\iii\in D(R):B(x,R)\cap J_{\iii}\cap E\neq \emptyset\}\\
			&\le 2R^{-{\eps}}\Big(\frac R r\Big)^{s^*(\eta)+2\eps}\sup_{x\in E}\#\{\iii\in D(R):B(x,R)\cap J_{\iii}\neq \emptyset\}\\
			&<2\Big(\frac R r\Big)^{s^*(\eta)+2\eps+\f{\eps}{\eta}}.
		\end{align*}
		Therefore, $h_E(2\eta)\le s^*(\eta)+2\eps+\f{\eps}{\eta}.$  Letting $\eps\to 0$ completes the proof.
	\end{proof}

\end{document}
